%==============================================================================
%  CAPÍTULO XX
%==============================================================================
\chapter{Gestión del encargo: mandato, honorarios y responsabilidad}
\label{cap:gestion-encargo}

\fichaCapitulo{La dimensión profesional del ejercicio concursal.}{Delimitación del
encargo, estructura de honorarios, deberes de información y prevención de la
responsabilidad.}

%==============================================================================
\section{Por qué el encargo concursal exige delimitación expresa}
%==============================================================================

El procedimiento concursal presenta una combinación poco frecuente: alta expectativa del
cliente, resultado parcialmente fuera del control del abogado y duración prolongada. Los
tres factores concurren para producir insatisfacción aun en procedimientos técnicamente
bien llevados.

El cliente llega buscando que sus deudas desaparezcan. El abogado sabe que algunas
sobrevivirán, que otras dependerán del voto de los acreedores y que el conjunto tomará
meses. Si esa distancia no se explicita al inicio y por escrito, se manifestará al final,
cuando ya no pueda gestionarse.

\begin{foco}[La regla práctica]
Todo lo que el cliente vaya a descubrir con desagrado durante el procedimiento debe
haberlo leído y firmado antes de que comience.
\end{foco}

%==============================================================================
\section{Contenido mínimo del mandato}
%==============================================================================

\begin{checklist}[Cláusulas que no deben faltar]
  \item \textbf{Objeto preciso.} Procedimiento específico que se tramitará, con exclusión
  expresa de los demás. La liquidación no incluye la defensa de las acciones revocatorias
  ni la representación en juicios laborales derivados.
  \item \textbf{Deudas excluidas.} Mención expresa de que el crédito con aval del Estado,
  el crédito solidario y las pensiones de alimentos probablemente no se extinguirán.
  \item \textbf{Carácter público.} Constancia de que el procedimiento se publica en el
  \boletin{} y es de acceso general.
  \item \textbf{Deber de información del cliente.} Obligación de declarar la totalidad de
  sus bienes y actos, con mención de las consecuencias del \art{169 A}.
  \item \textbf{Duración estimada.} Con advertencia de que los plazos son referenciales y
  dependen de la agenda del tribunal o de la \superir{}.
  \item \textbf{Honorarios.} Monto, oportunidad de pago y qué ocurre si el procedimiento
  se extiende o termina anticipadamente.
  \item \textbf{Gastos.} Certificados, receptores, publicaciones y su forma de provisión.
  \item \textbf{Término anticipado.} Condiciones de renuncia y de revocación, y efectos
  sobre los honorarios devengados.
\end{checklist}

\begin{alerta}[La cláusula sobre información del cliente]
No es una formalidad defensiva. Es la que permite renunciar al encargo con fundamento
cuando, avanzado el procedimiento, aparece el activo que el cliente había ocultado
deliberadamente. Sin ella, el abogado queda atrapado entre su deber de lealtad y su deber
de no colaborar con una declaración que sabe falsa.
\end{alerta}

%==============================================================================
\section{Estructura de honorarios}
%==============================================================================

\subsection{El problema de fondo}

El cliente concursal es, por definición, insolvente. La estructura de honorarios debe
resolver esa tensión sin comprometer la independencia profesional ni la viabilidad del
encargo.

\subsection{Modalidades y sus riesgos}

\begin{description}[leftmargin=0pt,style=nextline,font=\sffamily\bfseries]
  \item[Suma fija pagada al inicio] Es la más limpia y la más difícil de obtener. Evita
  el conflicto de interés y la incobrabilidad.

  \item[Suma fija en cuotas durante la tramitación] Frecuente en renegociación y
  liquidación simplificada. Riesgo: el cliente deja de pagar a mitad del procedimiento y
  el abogado queda obligado a continuar o a renunciar en un momento inoportuno.

  \item[Honorario por etapas] Vincula el pago a hitos verificables (presentación,
  resolución de admisibilidad, junta, término). Distribuye el riesgo y es fácil de
  explicar.

  \item[Honorario variable sobre el resultado] Debe tratarse con cautela. En materia
  concursal el «resultado» es difícil de definir y su vinculación al monto extinguido
  puede generar incentivos indeseables.
\end{description}

\begin{practica}[Provisión de gastos separada]
Los gastos ---certificados, publicaciones, receptores--- deben provisionarse por
separado y documentarse. Confundirlos con el honorario es una fuente habitual de disputa
al término del procedimiento.
\end{practica}

%==============================================================================
\section{Deberes de información durante la tramitación}
%==============================================================================

\begin{checklist}[Comunicaciones que deben quedar por escrito]
  \item Resolución que da inicio al procedimiento y sus efectos.
  \item Nominación del veedor o liquidador y qué implica para el cliente.
  \item Cada resolución que altere la posición del cliente.
  \item Verificaciones de crédito relevantes y decisión de objetarlas o no.
  \item Convocatoria a junta o audiencia, con explicación de lo que se resolverá.
  \item Resultado de la votación y sus consecuencias.
  \item Resolución de término y detalle de qué deudas se extinguieron y cuáles no.
\end{checklist}

\begin{foco}[El informe de cierre]
Al terminar el procedimiento conviene entregar un informe que consigne qué se obtuvo,
qué deudas subsisten y qué debe hacer el cliente en adelante. Cierra el encargo, previene
consultas posteriores y documenta el cumplimiento del deber de información.
\end{foco}

%==============================================================================
\section{Prevención de la responsabilidad profesional}
%==============================================================================

\errorfrecuente{Aceptar el encargo sin verificación documental previa}
{Se construye la estrategia sobre información falsa y el procedimiento fracasa por causas
imputables a la preparación.}
{Condicionar la aceptación a la entrega de los antecedentes del capítulo
\ref{cap:entrevista}.}

\errorfrecuente{No documentar las advertencias sobre deudas excluidas}
{Frente al reclamo del cliente no hay constancia de que fue informado.}
{Acta de diagnóstico firmada y cláusula expresa en el mandato.}

\errorfrecuente{Delegar el seguimiento del Boletín Concursal sin control}
{Se pierden plazos fatales cuya reapertura no existe.}
{Sistema de control con responsable identificado y revisión diaria documentada.}

\errorfrecuente{Continuar el encargo tras descubrir ocultamiento deliberado}
{Compromete al profesional en la construcción de una declaración que sabe inexacta.}
{Invocar la cláusula de información, exigir la rectificación y, en su defecto, renunciar
dejando constancia.}

\errorfrecuente{Prometer plazos como si fueran ciertos}
{La demora ordinaria del sistema se percibe como negligencia del abogado.}
{Informar rangos, explicar de qué dependen y actualizar periódicamente.}

\begin{alerta}[Sobre el conflicto de interés]
La asesoría simultánea a deudor y acreedores en un mismo concurso, o al deudor y a una
sociedad relacionada que resulte contraparte de una acción revocatoria, genera conflictos
que suelen advertirse tarde. Conviene mapear las relaciones del cliente antes de aceptar,
y no cuando el conflicto ya se materializó.
\end{alerta}
